\documentclass[11pt]{letter}
\usepackage[letterpaper,margin=0.9in]{geometry}
\usepackage[hidelinks]{hyperref}
\usepackage{parskip}
\pagenumbering{gobble}

\begin{document}

\begin{center}
{\Large \textbf{Aryan Miriyala}}\\
\href{mailto:aryanmiriyala@gmail.com}{aryanmiriyala@gmail.com} $|$ +1 419-315-0444 $|$
\href{https://www.linkedin.com/in/aryan-miriyala-7788a21b6/}{linkedin.com/in/aryan-miriyala} $|$
\href{https://github.com/aryanmiriyala}{github.com/aryanmiriyala}
\end{center}

\vspace{0.25in}

Dear Amazon Hiring Team,

I am applying for the Software Development Engineer - 2026 role because Amazon's SDE work is centered on ownership, customer impact, and building reliable systems at scale. The part of the role that stands out to me is the expectation that engineers own code from design through operation while solving real customer problems quickly and carefully.

My background fits that kind of early-career engineering path through internships and projects where I have built secure software, AWS-backed automation, database workflows, and AI-assisted tools. At AAIS, I automated a MySQL/SQL billing workflow with Python, PySpark, and AWS Glue over 20+ TB of production data for 700+ member companies, profiled 160+ source tables, and built AWS Lambda/S3 validation workflows. At SmartSolve, I built secure internal software with Next.js, TypeScript, PostgreSQL, SSO, auth middleware, and Docker-based development environments while using Codex and Claude Code for planning and review. I have also built distributed/data and systems projects with Kafka, Hadoop HDFS, Spark, C++, OpenMP, Docker, REST APIs, and AI evaluation workflows. I would bring Amazon a practical ownership mindset: learn quickly, communicate clearly, write maintainable code, debug carefully, and keep customer-facing reliability at the center of the work.

\vspace{0.15in}

Best regards,\\
Aryan Miriyala

\end{document}
